\section{The player--profile indexing condition}\label{sec:indexing}
Let $N=\{0,\ldots,n-1\}$, let $A$ be a common nonempty finite action set,
and let $u_i:A^n\to\R$. For $p\in S_n$, define
\begin{equation}\label{eq:profile-action}
 (P_p s)_j=s_{p^{-1}(j)}.
\end{equation}
Thus the action originally chosen by $i$ is used by $p(i)$ after relabeling,
and $P_qP_p=P_{qp}$. Under this convention the usual equivariance condition is
\begin{equation}\label{eq:correct-symmetry}
 u_i(s)=u_{p(i)}(P_p s)\qquad(i\in N,\ s\in A^n,\ p\in S_n).
\end{equation}
By contrast, pairing $p(i)$ with the non-inverse list of profile indices gives
\begin{equation}\label{eq:mismatched-symmetry}
 u_i(s)=u_{p(i)}(s_{p(0)},\ldots,s_{p(n-1)})
       =u_{p(i)}(P_{p^{-1}}s).
\end{equation}
These formulas coincide on transpositions. Their difference therefore cannot
be detected in two-player examples, or by checking only player transpositions.
A consistent right-action convention is equally valid; the issue is the
mismatched pairing, not a preference for left-action notation.

\begin{theorem}[Vester; Ham]\label{thm:indexing}
Suppose $n\ge3$. Condition~\eqref{eq:mismatched-symmetry} holds if and only if
\eqref{eq:correct-symmetry} holds and all payoff functions are identical:
$u_i(s)=u_j(s)$ for every $i,j,s$.
\end{theorem}
\begin{proof}
On each transposition, the mismatched condition is the correct condition.
The correctly preserving player permutations form a subgroup: successive
substitutions use $P_qP_p=P_{qp}$, and inversion follows by changing profiles.
Since transpositions generate $S_n$, the correct condition holds for all $p$.
Compare the mismatched condition at $p$ with the correct condition at $p^{-1}$.
They give
\[
 u_{p(i)}(P_{p^{-1}}s)=u_{p^{-1}(i)}(P_{p^{-1}}s).
\]
The profile map is bijective. Given distinct $a,b$, choose a third player $k$
and take the cycle $p=(k\ a\ b)$ and $i=k$; then $p(i)=a$ and $p^{-1}(i)=b$.
Hence $u_a=u_b$. Conversely apply the correct condition at $p^{-1}$ and
use equality of the recipient payoff functions.
\end{proof}
The algebraic proof uses no finiteness of $A$, although our game and equilibrium
statements remain finite. The common-payoff consequence was publicly proved
by \citet[Appendix B, Theorem 32]{vester2012}; the full equivalence appears in
\citet[Theorem 3.18]{ham2013} and \citet[Theorem 3.21]{ham2021}. See also
\citet[Lemma 2.14 and Proposition 2.15]{tohme2019}, using its identified author
preprint for the result numbering. We do not claim a new discovery of this
indexing theorem.

\begin{example}\label{ex:indexing}
For $n\ge3$ and $A=\{0,1\}$, the game $u_i(s)=s_i$ satisfies
\eqref{eq:correct-symmetry} but not~\eqref{eq:mismatched-symmetry}. For example,
with $p=(0\ 1\ 2)$, $i=0$ and $s=(1,0,0,\ldots,0)$, the latter compares
$1$ with $0$. For $n=2$ both formulas hold in this same game despite the
payoff functions being distinct. Conversely identical payoff functions alone
do not imply either condition: $u_i(s)=s_0$ for every $i$ is not equivariant.
The mismatched condition does not force constant payoffs across profiles:
common payoff $\sum_j s_j$ satisfies it. If the game is zero sum and $n\ge3$,
Theorem~\ref{thm:indexing} instead forces every payoff to be zero at every
profile, since all $n$ payoffs are equal and sum to zero.
\end{example}

\paragraph{Source and scope of the correction.}
The mismatched formula is attributed to Definition 7 of
\citet{dm1986} in the inspected proofs of \citet{vester2012,ham2021} and the
author draft of \citet{plan2017}. Our mathematical assertion concerns the
explicit formula~\eqref{eq:mismatched-symmetry}. It neither asserts an error
in all permutation conventions nor claims to invalidate the discontinuous-
game equilibrium theorems of \citet{dm1986}. A source-level assessment of those
arguments is a separate question. Ham's 2011-dated thesis identifies the
indexing problem \citep[Subsection~5.8]{hamthesis}; its public arXiv posting is 2019. The 2012 and 2013 public
proofs, rather than that thesis date, determine the attribution of the theorem.
